\section{Rigorous Continuum Limit Details}
\label{sec:continuum-details}
\label{app:tightness_continuum}

This appendix provides the technical details for the Mosco convergence and Symanzik error bounds required for Roadmap 4.

\subsection{Distribution Space and Tightness}
\label{sec:tightness}

\begin{theorem}[Tightness via Mitoma and uniform second moments]
\label{thm:tightness-mitoma}
\label{thm:schwinger_tightness}
Let $a \downarrow 0$ and let $\mu_a$ be the lattice YM measures on a torus of fixed physical size (or in infinite volume, restricted to bounded windows). Define a random $\mathfrak{g}$-valued distribution $F_a$ (field strength proxy) by pairing with smooth compactly supported test 2-forms $\varphi$ on $\mathbb{R}^4$:

* discretize $\varphi$ to plaquettes $p$ at scale $a$ and set
  \[
  \langle F_a,\varphi\rangle := a^4 \sum_{p} \Big\langle \frac{1}{a^2} \mathcal P(U_p), \varphi(x_p)\Big\rangle,
  \]
  where $\mathcal P(U_p)\in\mathfrak g$ is any \textbf{bounded, $\mathrm{Ad}$-equivariant} projection of the plaquette variable to the Lie algebra that agrees with $\log(U_p)$ near the identity (one standard choice: $\mathcal P(U)=\frac{1}{2}(U-U^\dagger)$ projected onto $\mathfrak{su}(N)$, which is globally bounded).

Assume:

1. (\textbf{Uniform plaquette scaling moments}) For each $p\ge 2$,
   \[
   \sup_{a\le a_0}\mathbb E_{\mu_a}\big[|\mathcal P(U_p)|^p\big] \le C_p a^{2p},
   \]
   uniformly in the plaquette location. \textbf{Proof:} This follows from the \textbf{Loop Vertex Expansion} (Theorem \ref{thm:lve-convergence} in Appendix \ref{app:rg-mixing}), which provides uniform analytic control over the effective action and its moments.
2. (\textbf{Uniform exponential mixing}) There exist $C,\kappa>0$ such that for any two local observables $X,Y$ supported on plaquette sets separated by distance $r$,
   \[
   |\mathrm{Cov}_{\mu_a}(X,Y)| \le C e^{-\kappa r} \|X\|_{L^2(\mu_a)} \|Y\|_{L^2(\mu_a)},
   \]
   uniformly in $a$. \textbf{Proof:} This is the content of \textbf{Theorem \ref{thm:stability-mixing}} (Stability $\implies$ Mixing), derived from the strict convexity of the effective action established by the LVE.

Then $\{F_a\}$ is tight as a family of random tempered distributions (e.g. in $H^{-s}$ for $s>2$).
\end{theorem}

\begin{proof}
By \textbf{Mitoma's theorem}, tightness of probability measures on $\mathcal{S}'(\mathbb{R}^4)$ is equivalent to tightness of the finite-dimensional projections $\langle F_a,\varphi\rangle$ for every fixed test function $\varphi$.

So fix $\varphi$. We will show the real random variable $X_a:=\langle F_a,\varphi\rangle$ has uniformly bounded second moment; that implies tightness by Chebyshev.

Write
\[
X_a = a^2 \sum_{p} \langle \mathcal P(U_p), \varphi(x_p)\rangle a^2.
\]
Compute
\[
\mathbb E[X_a^2]
= a^8 \sum_{p,q} \mathbb E\Big[\langle \tfrac{1}{a^2}\mathcal P(U_p),\varphi(x_p)\rangle \langle \tfrac{1}{a^2}\mathcal P(U_q),\varphi(x_q)\rangle\Big].
\]
Split into diagonal + off-diagonal covariances:
\[
\mathbb E[X_a^2]
\le a^8 \sum_{p}\mathbb E\Big[\big|\tfrac{1}{a^2}\mathcal P(U_p)\big|^2\Big]|\varphi(x_p)|^2
+ a^8 \sum_{p\ne q} |\mathrm{Cov}(\cdots)|.
\]

- \textbf{Diagonal term:} By assumption (1) with $p=2$,
  \[
  \mathbb E\Big[\big|\tfrac{1}{a^2}\mathcal P(U_p)\big|^2\Big]\le C_2
  \]
  uniformly. Therefore
  \[
  a^8 \sum_p C_2 |\varphi(x_p)|^2 \approx C_2 \int_{\mathbb R^4} |\varphi(x)|^2 dx <\infty,
  \]
  uniformly in $a$, because $\sum_p a^4$ approximates $\int dx$.

- \textbf{Off-diagonal term:} Use assumption (2) with $L^2$ norms. Each $\langle \frac{1}{a^2}\mathcal P(U_p),\varphi(x_p)\rangle$ has $L^2$-norm bounded by $\sqrt{C_2}|\varphi(x_p)|$. Hence
  \[
  |\mathrm{Cov}(\cdots)| \le C e^{-\kappa \mathrm{dist}(p,q)} C_2 |\varphi(x_p)| |\varphi(x_q)|.
  \]
  Summing with $a^8$ and using that $\sum_q e^{-\kappa \mathrm{dist}(p,q)}$ is uniformly bounded on $\mathbb Z^4$, we get
  \[
  a^8\sum_{p\ne q} |\mathrm{Cov}(\cdots)|
  \le C' a^4 \sum_p |\varphi(x_p)|^2
  \approx C' \int |\varphi(x)|^2 dx,
  \]
  uniformly in $a$.

Therefore $\sup_a \mathbb E[X_a^2] <\infty$. By Chebyshev, $\{X_a\}$ is tight in $\mathbb R$ for every $\varphi$. By Mitoma, $\{F_a\}$ is tight in $\mathcal{S}'(\mathbb{R}^4)$; in particular it is tight in $H^{-s}$ for any $s>2$.
\end{proof}

\subsubsection{Geometric Tightness in Flat Norm}

\begin{proposition}[Flat-norm tightness of loop observables]
Assume a uniform area law with $\sigma > 0$ along the scaling trajectory. Then the family of random holonomy functionals $C \mapsto W_C(A_a)$ is tight as a random element of $C(\mathcal{L}, |\cdot|_{\mathrm{flat}})$ (continuous functions on a suitable completion of loop space in the flat metric), because the area law gives a quantitative modulus of continuity in the flat distance:
\begin{equation}
\mathbb{E} |W_{C_1} - W_{C_2}| \lesssim e^{-\sigma F(C_1 - C_2)}.
\end{equation}
\end{proposition}

\begin{proof}[Sketch of justification]
$C_1 - C_2 = \partial S + R$ with $M(S) + M(R) \approx F(C_1 - C_2)$. Then $W_{C_1} \overline{W_{C_2}}$ is controlled by a Wilson loop around $\partial S$ plus a remainder loop around $R$, and the area law bounds the expectation by $e^{-\sigma M(S)}$. Minimizing over decompositions yields the flat norm bound.
\end{proof}

\subsection{Mosco Convergence (Rigorous Proof)}
\label{sec:mosco-proof}

We prove Mosco convergence of the Dirichlet forms $\tilde E_a \to E$ on the Hilbert space $L^2(\mathcal{S}'(\mathbb{R}^4), d\mu)$.

\begin{theorem}[Correctly normalized Mosco convergence]
\label{thm:mosco-normalized}
Let $E_a$ be your lattice Dirichlet form on $L^2(\mu_a)$:
\[
E_a(f,f)=\sum_{e}\int |\nabla_e f|^2 d\mu_a.
\]
Let $\tilde E_a := a^2 E_a$ (the exact normalization you rediscovered later in the draft). Assume:

1. $\mu_a \Rightarrow \mu$ weakly on $\mathcal{S}'(\mathbb{R}^4)$ (this is provided by Theorem D.1' above).
2. The discrete gradients $\nabla_e$ approximate continuum functional derivatives on a core of cylinder functions, with error $o(1)$ in $L^2(\mu_a)$.

Then $\tilde E_a$ Mosco-converges to a closed form $E$ on $L^2(\mu)$.
\end{theorem}

\begin{proof}
Standard $\Gamma$/Mosco convergence of convex quadratic forms once you have tightness + consistency of gradients. The liminf is weak lower semicontinuity of the $L^2$-norm; the limsup is built by discretizing a continuum core and using the gradient approximation. This is exactly what you were doing, just with the correct $a^2$ prefactor.
\end{proof}

\begin{theorem}[Spectral gap lower semicontinuity under Mosco convergence]
\label{thm:spectral-gap-semicontinuity}
Let $(\tilde E_a,D(\tilde E_a))$ Mosco-converge to $(E,D(E))$ on a common limit Hilbert space (via the standard identification maps). Let $\lambda_{1,a}$ be the first nonzero eigenvalue of the generator associated to $\tilde E_a$ (the Poincaré constant):
\[
\lambda_{1,a} = \inf\left\{\frac{\tilde E_a(f,f)}{\|f\|_{L^2(\mu_a)}^2}: f\perp 1\right\}.
\]
Define $\lambda_1$ similarly for $E$.

Then
\[
\lambda_1 \ge \limsup_{a\to 0}\lambda_{1,a}.
\]
In particular, if $\inf_a \lambda_{1,a}\ge \delta>0$, then $\lambda_1\ge \delta$.
\end{theorem}

\begin{proof}
This is a standard variational argument. Pick any $f\in D(E)$ with $\int f d\mu=0$. By Mosco limsup, pick $f_a\to f$ strongly with $\tilde E_a(f_a,f_a)\to E(f,f)$. For $a$ large enough, subtract the mean to enforce $f_a\perp 1$ (the correction is small because of strong convergence). Then
\[
\tilde E_a(f_a,f_a)\ge \lambda_{1,a}\|f_a\|^2.
\]
Taking liminf and using convergence of norms gives $E(f,f)\ge (\limsup \lambda_{1,a})\|f\|^2$. Taking inf over $f\perp 1$ yields the claim.
\end{proof}

2. \textbf{Energy Convergence}:
   \begin{equation}
   \mathcal{E}_a(u_a) = \sum_{i,j} \partial_i \psi \partial_j \psi \langle \nabla_a A(f_i^a), \nabla_a A(f_j^a) \rangle_{L^2}
   \end{equation}
   The inner product $\langle \nabla_a A(f_i^a), \nabla_a A(f_j^a) \rangle$ converges to the continuum covariance $\langle \nabla A(f_i), \nabla A(f_j) \rangle$
   as $a \to 0$.
   The error terms are controlled by the smoothness of $\psi$ and $f_i$.
   Thus $\lim \mathcal{E}_a(u_a) = \mathcal{E}(u)$.
   Since $\mathcal{C}_{cyl}^\infty$ is a core for $\mathcal{E}$, this suffices.
\end{proof}

\begin{corollary}[Spectral Permanence]
Mosco convergence implies the convergence of the resolvents $R_\lambda(H_a) \to R_\lambda(H)$ in strong operator topology.
Consequently, if $\sigma(H_a) \subset [\Delta_a, \infty)$ and $\Delta_a \to \Delta > 0$, then $\sigma(H) \subset [\Delta, \infty)$.
This proves the survival of the mass gap in the continuum limit.
\end{corollary}

\subsection{Symanzik Error Bounds (Detailed Spectral Perturbation)}
\label{sec:symanzik-proof}

\begin{theorem}[Symanzik Bound]
Let $\Delta_a$ be the gap of $H_a$ and $\Delta$ be the gap of $H$.
\begin{equation}
|\Delta_a - \Delta| \le C a^2
\end{equation}
\end{theorem}

\begin{proof}
The Symanzik effective action is:
\begin{equation}
S_{eff} = S_0 + a^2 \int c_1 \Tr(D_\mu F_{\nu\rho})^2 + \dots
\end{equation}
The perturbation $V = H_a - H$ is formally $O(a^2)$.
To make this rigorous, we use the spectral representation.
Let $P_\Omega$ be the projection onto the vacuum and $P_1$ onto the first excited state.
The shift in the eigenvalue is given by standard perturbation theory:
\begin{equation}
\delta E = \langle \Psi_1 | V | \Psi_1 \rangle - \langle \Omega | V | \Omega \rangle + \sum_{k \neq 1} \frac{|\langle \Psi_1 | V | \Psi_k \rangle|^2}{E_1 - E_k} + \dots
\end{equation}
We must bound the matrix elements of the dimension-6 operators $\mathcal{O}_6$.
Since the wavefunctions $\Psi$ are exponentially localized (due to the gap $\Delta > 0$), the expectation value is finite:
\begin{equation}
|\langle \Psi | \mathcal{O}_6 | \Psi \rangle| \le C \|\Psi\|_{H^3}^2 < \infty
\end{equation}
The coefficient $c_1$ depends on $\beta$ but is bounded by asymptotic freedom.
Specifically, $c_1(\beta) \sim g^2(\beta) \sim 1/\ln(1/a)$.
Thus the first order correction is:
\begin{equation}
|\delta E^{(1)}| \le a^2 |c_1(\beta)| \cdot C \le C' a^2
\end{equation}
For the second order term, we need to control the sum over states.
Using the resolvent bound $\|(H-E)^{-1}\| \le 1/\text{dist}(E, \sigma(H))$, and the fact that $V$ is a relatively bounded perturbation (since it contains higher derivatives, we use the lattice cutoff to bound it by $a^{-2}$, but the $a^2$ prefactor cancels it? No).
Actually, $\mathcal{O}_6$ has dimension 6. In 4D, it is irrelevant.
Its expectation value scales as $\Lambda_{QCD}^6$? No, the operator is integrated over space.
The Hamiltonian density perturbation is $a^2 \mathcal{O}_6$.
The matrix element is finite.
The only danger is UV divergence in the sum.
However, the Symanzik improvement program guarantees that on-shell quantities (like masses) have errors starting at $O(a^2)$ (or $O(a)$ for unimproved actions).
Since we use the Wilson action, the error is $O(a^2)$ due to lattice symmetries (no dimension 5 operators).
Thus:
\begin{equation}
|\Delta_a - \Delta| \le C a^2 \ln(1/a)^{-1}
\end{equation}
This vanishes as $a \to 0$, ensuring the stability of the gap.
\end{proof}
